\documentclass[10pt,a4paper]{article}
\usepackage[a4paper,margin=14mm,top=12mm,bottom=14mm]{geometry}
\usepackage{fontspec}
\setmainfont{Liberation Sans}
\usepackage{polyglossia}
\setdefaultlanguage{french}
\usepackage{amsmath}
\usepackage{xcolor}
\usepackage{graphicx}
\usepackage[most]{tcolorbox}
\usepackage{tabularx,booktabs,array,multirow}
\usepackage{enumitem}
\usepackage{fancyhdr}
\usepackage{lastpage}
\usepackage[hidelinks]{hyperref}
\definecolor{uimm}{HTML}{005677}
\definecolor{cyan}{HTML}{13A2DD}
\definecolor{fond}{HTML}{E6F2FA}
\definecolor{ok}{HTML}{2E7D32}
\definecolor{ko}{HTML}{B3261E}
\definecolor{violet}{HTML}{7B1FA2}
\setlength{\parindent}{0pt}
\setlength{\parskip}{3pt}
\setlist{nosep}
\setlist[itemize]{leftmargin=13pt,label=\textcolor{cyan}{\textbullet}}
\newcommand{\titre}[1]{\vspace{5pt}{\color{uimm}\large\bfseries #1}\par\vspace{-3pt}{\color{cyan}\rule{\linewidth}{1.2pt}}\par\vspace{0pt}}
\tcbset{boite/.style={colback=fond,colframe=uimm,boxrule=0.6pt,arc=2pt,left=6pt,right=6pt,top=4pt,bottom=4pt},
        defi/.style={colback=violet!6,colframe=violet,boxrule=0.8pt,arc=2pt,left=6pt,right=6pt,top=4pt,bottom=4pt,
                     title={\bfseries Défi},coltitle=white,colbacktitle=violet,fonttitle=\small}}
\newcommand{\bandeau}[3]{%
\begin{tcolorbox}[colback=uimm,colframe=uimm,arc=2pt,left=4pt,right=8pt,top=4pt,bottom=4pt]
\begin{minipage}[c]{18mm}\setlength{\fboxsep}{1.2pt}\fcolorbox{white}{white}{\includegraphics[height=16.5mm]{logo_uimm.jpg}}\end{minipage}\hspace{4mm}%
\begin{minipage}[c]{\dimexpr\linewidth-22mm\relax}\color{white}
{\small #1}\\[1pt]{\LARGE\bfseries #2}\\[2pt]{\large #3}\hfill{\small CQPM Technicien qualité}
\end{minipage}
\end{tcolorbox}}
\newcommand{\reponse}[1][1]{\par\vspace{2pt}{\color{gray!60}\foreach \i in {1,...,#1}{\rule{\linewidth}{0.3pt}\par\vspace{9pt}}}}
\usepackage{pgffor}
\pagestyle{fancy}\fancyhf{}
\renewcommand{\headrulewidth}{0pt}
\newcommand{\pied}[1]{\fancyfoot[L]{\footnotesize\color{uimm}#1}\fancyfoot[R]{\footnotesize\color{uimm}Page \thepage{} sur \pageref{LastPage}}}

\usepackage{tikz}
\usepackage{longtable}
\renewcommand{\bandeau}[3]{%
\begin{tcolorbox}[colback=uimm,colframe=uimm,arc=2pt,left=4pt,right=8pt,top=4pt,bottom=4pt]
\begin{minipage}[c]{18mm}\setlength{\fboxsep}{1.2pt}\fcolorbox{white}{white}{\includegraphics[height=16.5mm]{logo_uimm.jpg}}\end{minipage}\hspace{4mm}%
\begin{minipage}[c]{\dimexpr\linewidth-22mm\relax}\color{white}
{\small #1}\\[1pt]{\LARGE\bfseries #2}\\[2pt]{\large #3}
\end{minipage}
\end{tcolorbox}}
% en-tête d'une fiche : code, titre, question, bloc, durée, organisation, niveau
\newcommand{\entete}[7]{%
\begin{tcolorbox}[colback=uimm,colframe=uimm,arc=2pt,left=4pt,right=6pt,top=4pt,bottom=4pt]
\begin{minipage}[c]{18mm}\setlength{\fboxsep}{1.2pt}\fcolorbox{white}{white}{\includegraphics[height=16.5mm]{logo_uimm.jpg}}\end{minipage}\hspace{3mm}%
\begin{minipage}[c]{\dimexpr\linewidth-58mm\relax}\color{white}\raggedright
{\small Base SPC · #4}\\[1pt]{\Large\bfseries #1 · #2}\\[2pt]{\itshape #3}
\end{minipage}\hfill
\begin{minipage}[c]{33mm}\color{white}\small\raggedleft
Durée : \textbf{#5}\\ \textbf{#6}\\ Niveau #7 sur 3
\end{minipage}
\end{tcolorbox}}
\newcommand{\tagE}{{\setlength{\fboxsep}{1.5pt}\colorbox{cyan}{\color{white}\scriptsize\bfseries Ellistat}}}
\newcommand{\tagW}{{\setlength{\fboxsep}{1.5pt}\colorbox{uimm}{\color{white}\scriptsize\bfseries Page web}}}
\tcbset{pred/.style={colback=orange!6,colframe=orange!80!black,boxrule=0.6pt,arc=2pt,left=6pt,right=6pt,top=3pt,bottom=1pt,
                    title={\bfseries Je prédis, avant de commencer},coltitle=white,colbacktitle=orange!80!black,fonttitle=\small},
        retenir/.style={colback=fond,colframe=uimm,boxrule=0.6pt,arc=2pt,left=6pt,right=6pt,top=3pt,bottom=3pt,
                    title={\bfseries À retenir},coltitle=white,colbacktitle=uimm,fonttitle=\small},
        plus/.style={colback=violet!6,colframe=violet,boxrule=0.6pt,arc=2pt,left=6pt,right=6pt,top=3pt,bottom=3pt,
                    title={\bfseries Pour aller plus loin},coltitle=white,colbacktitle=violet,fonttitle=\small},
        act/.style={colback=uimm!10,colframe=uimm,boxrule=0.5pt,arc=2pt,left=5pt,right=5pt,top=2pt,bottom=2pt,before skip=7pt,after skip=3pt},
        verif/.style={colback=orange!6,colframe=orange!80!black,boxrule=0.6pt,arc=2pt,left=6pt,right=6pt,top=3pt,bottom=3pt,
                    title={\bfseries À vérifier sur poste avant la première séance},coltitle=white,colbacktitle=orange!80!black,fonttitle=\small}}
\newcommand{\fichier}[1]{\texttt{\small #1}}
\newcommand{\menu}[1]{\textsf{\textbf{#1}}}
\newcommand{\champ}[1]{\rule{#1}{0.3pt}}

\pied{Base SPC · fiche M6 · La cale dans la main}
\begin{document}
\thispagestyle{fancy}
\entete{M6}{La cale dans la main}{La chaleur des doigts se voit-elle au micromètre ?}{Le moyen de mesure}{25 min}{en binôme}{2}
\begin{tcolorbox}[boite]\textbf{Le but.} Voir l'effet de la température sur une mesure au millième.\par\textbf{Il vous faut.} Un micromètre 0-25 mm · Une cale étalon de 25 mm (ou la plus longue disponible) · Un chronomètre\end{tcolorbox}
\begin{tcolorbox}[pred]Une cale d'acier de 25 mm tenue 3 minutes dans la main fermée s'allonge-t-elle ? De combien : 0,1 µm, 3 µm, 30 µm ?\reponse\end{tcolorbox}
\titre{À faire}
\begin{enumerate}[label=\textcolor{cyan}{\textbf{\arabic*.}},leftmargin=15pt,itemsep=1.5pt]
\item Mesurer 3 fois la cale à la température de la salle. Noter la moyenne : c'est la référence.
\item Garder la cale 3 minutes dans la main fermée.
\item La poser et la mesurer aussitôt, puis toutes les 30 secondes pendant 5 minutes, sans la reprendre en main.
\item Tracer les points sur le graphique : temps en abscisse, mesure en ordonnée.
\end{enumerate}
\vspace{3pt}Référence (moyenne de 3 mesures à la température de la salle) : \champ{22mm} mm\par\vspace{3pt}
{\renewcommand{\arraystretch}{1.5}\small\begin{tabularx}{\linewidth}{|l|*{11}{>{\centering\arraybackslash}X|}}\hline
\textbf{Temps (s)} & 0 & 30 & 60 & 90 & 120 & 150 & 180 & 210 & 240 & 270 & 300\\\hline
\textbf{Mesure (mm)} &  &  &  &  &  &  &  &  &  &  & \\\hline\end{tabularx}}\par\vspace{4pt}
\begin{center}\begin{tikzpicture}[x=1.45cm,y=0.5cm]\draw[gray!45,very thin] (0,0) grid (10,6);\draw[thick] (0,0)--(10,0);\draw[thick] (0,0)--(0,6);\node[below,font=\scriptsize] at (0,0) {0};\node[below,font=\scriptsize] at (1,0) {30};\node[below,font=\scriptsize] at (2,0) {60};\node[below,font=\scriptsize] at (3,0) {90};\node[below,font=\scriptsize] at (4,0) {120};\node[below,font=\scriptsize] at (5,0) {150};\node[below,font=\scriptsize] at (6,0) {180};\node[below,font=\scriptsize] at (7,0) {210};\node[below,font=\scriptsize] at (8,0) {240};\node[below,font=\scriptsize] at (9,0) {270};\node[below,font=\scriptsize] at (10,0) {300};\node[below,font=\scriptsize] at (5,-0.55) {temps (s)};\node[rotate=90,font=\scriptsize] at (-0.55,3) {mesure : un carreau = 0,001 mm};\end{tikzpicture}\end{center}
\newpage
\titre{J'observe, je décide}
\begin{enumerate}[label=\textcolor{cyan}{\textbf{Q\arabic*.}},leftmargin=22pt]
\item De combien la cale s'est-elle allongée ? En combien de temps revient-elle à sa longueur ?\reponse[2]
\item L'acier s'allonge d'environ 11,5 µm par mètre et par degré. Pour 25 mm et 10 °C de plus, calculez l'allongement attendu.\reponse[2]
\item Quelles précautions en tirer au poste de contrôle ?\reponse[2]
\end{enumerate}
\begin{tcolorbox}[retenir]Au millième, la température compte : pièces et instruments à la même température, et on ne garde pas la pièce dans la main.\end{tcolorbox}
\begin{tcolorbox}[plus]Saisir les mesures dans Ellistat (une colonne, Répétition = 1) et regarder la carte : la dérive se voit-elle ?\end{tcolorbox}
\vfill{\footnotesize\color{uimm}\textbf{Ressources de la base :} Données de secours.}
\end{document}